% !TeX root = ../thuthesis-example.tex
\chapter{多模态方言声学知识图谱的构建}

本项目所有上层算法与框架的基石，是一个内容丰富、结构严谨、且专门面向方言学习的“多模态方言声学知识图谱”。它并非一个简单的词汇库，而是一个深度融合了语言学理论、声学特征与多媒体信息的认知中枢。构建这样一个高质量的知识图谱，是确保整个平台智能化、精准化、可解释性的前提。其构建过程遵循了从顶层理论设计到底层数据融合的系统化方法，涵盖了模式设计、多源数据融合、知识抽取与计算化表达等关键环节。

\section{知识图谱顶层模式（Schema）设计}
一个科学、完备的顶层模式是知识图谱的骨架。在构建初期，本项目首先对需要表达的方言知识进行了系统性的梳理，并设计了一套能够精确描述其复杂特征的Schema。

\subsection{核心实体与本体论基础}
本项目定义了多个核心实体类别，包括 \texttt{词汇 (Lexical Entry)}、\texttt{音节 (Syllable)}、\texttt{音系规则 (Phonological Rule)}、\texttt{发音人 (Speaker)} 以及 \texttt{多模态资产 (Multimodal Asset)}。为了确保模型的规范性和互操作性，本项目的 Schema 设计借鉴了万维网联盟（W3C）推荐的词汇本体模型 \textbf{OntoLex-Lemon}。该模型提供了一个连接词汇信息与本体概念的标准框架，其核心组件包括：
\begin{itemize}
    \item \texttt{Lexical Entry}: 词汇条目，代表一个独立的词。
    \item \texttt{Form}: 词形，一个词汇条目的具体书写或发音形式。
    \item \texttt{Sense}: 词义，将一个词汇条目链接到一个本体概念。
    \item \texttt{Concept}: 本体中的概念，代表词义的实际所指。
\end{itemize}
采用 OntoLex-Lemon 作为基础，使得本项目的知识图谱能够与更广泛的语言学数据资源（如 WordNet）进行链接与对齐。

\subsection{音系与超音段特征的扩展建模}
标准 OntoLex-Lemon 模型主要关注词汇和形态，对于音系学的细粒度信息，特别是方言中至关重要的超音段特征（Suprasegmental Features）涉及不足。因此，本项目对其进行了扩展：
\begin{enumerate}
    \item \textbf{音系表征}：本项目通过 OntoLex 的 \texttt{phon} (Phonology) 模块扩展 \texttt{Form} 实体，为其增加属性，用以存储国际音标（IPA）转写、音素序列等音系表征。
    \item \textbf{超音段特征建模}：对于声调、语调、重音和节奏等超音段特征，本项目将其模型化为 \texttt{多模态资产} 的一部分或 \texttt{音节} 实体的直接属性。例如，一个声调可以被定义为具有基频 $F_0$ 轮廓（Contour）、时长（Duration）等属性的实体。这使得一个“词汇”实体通过 \texttt{hasPronunciation} 关系连接到一个“音节”实体，该音节不仅包含音素信息，还关联着具体的声调模型，如图 \ref{fig:f0_contour} 所示。
    \item \textbf{多模态资产的深度定义}：其中，“多模态资产”实体不仅包含了常规的音频（Audio）、口型视频（Lip-sync Video）等，更将 \textbf{语图（Spectrogram）} 视为一种核心的、结构化的声学特征资产。此外，本项目进一步将其他关键声学特征也纳入其中：
        \begin{itemize}
            \item \textbf{梅尔频率倒谱系数 (MFCCs)}: 一个包含例如39个系数的向量，高效地表征了音色的核心特征。
            \item \textbf{共振峰 (Formants)}: $F_1, F_2, F_3$ 等共振峰频率，直接反映了发音时声道的状态，是元音区分的关键。
            \item \textbf{基频 ($F_0$)}: 反映声带振动频率，是声调和语调的主要物理载体。
        \end{itemize}
\end{enumerate}
\begin{figure}[H]
    \centering
    \includegraphics[scale=0.7]{八.pdf} 
    \caption{苏州话“八”的示意图}
    \label{fig:f0_contour}
\end{figure}

这些实体之间通过丰富的关系类型进行连接，例如，一个“词汇”实体通过 \texttt{hasPronunciation} 关系连接到一个“音节”实体，并通过 \texttt{hasAudioExample} 关系连接到一个“多模态资产”（音频文件）；该音频文件则进一步通过 \texttt{hasSpectrogram}, \texttt{hasMfccVector}, \texttt{hasF0Contour} 等关系连接到其对应的具体特征资产。


\section{多源异构数据采集与对齐}
知识图谱的内容来源于多个渠道的数据。语言学数据的采集主要依赖于权威的方言词典、地方志以及现代语言学研究专著。声学及视觉数据，则通过对目标方言的母语发音人，在专业录音环境下进行系统性的语料录制而获取。

\subsection{声学-文本强制对齐}
数据处理环节的核心难点在于多模态信息的对齐。本项目采用音素级别的 \textbf{强制对齐（Forced Alignment）} 技术，将录制的音频流与文本语料进行精确的时间戳匹配。考虑到方言往往是低资源语言，缺乏现成的声学模型和发音词典，本项目采用了以下策略：
\begin{itemize}
    \item \textbf{主流工具应用}：使用如 \textbf{Montreal Forced Aligner (MFA)} 等基于 Kaldi 工具包的先进对齐器。
    \item \textbf{低资源策略}：在目标方言缺乏训练数据时，采用“跨语言对齐”（Cross-language Alignment）策略，即利用一个声学特征相近且资源丰富语言（如普通话）的预训练模型进行初步对齐，再在少量方言数据上进行微调。
\end{itemize}
对齐的输出是每个音素、音节和词汇在音频文件中的精确起止时间，这是后续所有特征提取的基础。

\subsection{多模态特征提取与结构化}
在对齐完成后，本项目进一步对所有标准发音的音频和视频片段进行处理，提取出关键的声学与视觉特征并进行结构化存储。
\begin{enumerate}
    \item \textbf{声学“指纹”生成}: 为每个音节或词汇的标准发音生成并保存高分辨率的语图。这些语图作为一种可视化的声学“指纹”，被直接链接到对应的词汇或音节节点上，它们不仅是AI模型进行精准比对的依据，也可以作为一种直观的教具呈现给学习者。
    \item \textbf{数值特征提取}: 批量化提取前述的 MFCC 向量、共振峰轨迹和 $F_0$ 曲线。这些数值化特征被存储为独立的、可查询的文件（如 .csv 或 .npy），并在知识图谱中建立链接。
    \item \textbf{视觉特征提取}: 对口型视频进行处理，利用面部关键点检测模型定位唇部轮廓，提取出随时间变化的几何特征序列。
\end{enumerate}

\section{知识抽取、融合与存储实现}
知识图谱的构建是一个系统性工程，总体上可划分为知识抽取、知识学习（或融合）和知识评估三个阶段，如图 \ref{fig:知识图谱构建流程图} 所示。

\begin{figure}[H]
    \centering
    % 此处应放置您下载的知识图谱构建流程图
    \includegraphics[scale=0.23]{知识库.pdf}
    \caption{知识图谱构建流程图}
    \label{fig:知识图谱构建流程图}
\end{figure}

\subsection{自动化知识抽取}
在完成数据采集与对齐的基础上，本项目通过自动化的知识抽取脚本，将这些多源信息转化为知识图谱中的“主-谓-宾”三元组（Triples）。
\begin{itemize}
    \item \textbf{从结构化/半结构化数据抽取}: 对于方言词典等格式相对固定的数据源，本项目编写定制化的解析器，直接抽取 \texttt{(词汇A, 包含音节, 音节B)}、\texttt{(词汇A, 词义为, 概念C)} 等三元组。
    \item \textbf{从非结构化文本抽取}: 对于语言学专著等纯文本，本项目采用自然语言处理（NLP）流水线，包括命名实体识别（NER）以识别语言学术语，以及关系抽取（RE）模型来发现实体间的关系。常用的模型如基于BiLSTM-CRF的NER模型（如图 \ref{fig:ner_model} 所示 ）和基于卷积神经网络的抽取(RE)模型。
    \item \textbf{实体链接与消歧}: 方言中普遍存在同音词、多音字以及不同文献中转写不一的问题。本项目引入实体链接（Entity Linking）技术，如利用基于图的算法（类似 Babelfy 的思想），通过上下文和图谱的局部稠密度来判断一个词的提及（mention）应该链接到哪个唯一的词汇实体，从而确保知识的准确性。
\end{itemize}
例如，从词典中抽取 \texttt{(词汇: “食饭”, 包含音节, “sik6”)}，从对齐后的音频中标注 \texttt{(词汇: “食饭”, 拥有标准发音, 音频文件: “shifan\_male\_01.wav”)}，并基于声学分析结果建立 \texttt{(音频文件: “shifan\_male\_01.wav”, 生成语图, 语图文件: “shifan\_male\_01.png”)} 这样的关系。

\begin{figure}[H]
    \centering
    \includegraphics[width=\textwidth]{知识图谱使用.pdf}
    \caption{一种用于方言声学知识图谱的NER模型}
    \label{fig:ner_model}
\end{figure}


\subsection{知识融合与向量化存储}
所有抽取出的三元组经过清洗和实体对齐后，被融合进统一的知识库中。为实现高效的查询与管理，构建完成的知识图谱最终被存储在高性能的图数据库中。

同时，为了便于下游的机器学习模型直接调用，本项目利用知识图谱嵌入（Knowledge Graph Embedding）算法，将整个图谱中的实体和关系映射到低维、稠密的向量空间中。本项目考察了从经典的翻译模型（如 \textbf{TransE}、\textbf{RotatE}）到能够处理多模态信息的更先进模型的演进：
\begin{itemize}
    \item \textbf{传统模型}: 将知识图谱的结构信息编码为向量。
    \item \textbf{多模态嵌入模型}: 如 \textbf{IKRL} 或 \textbf{TransAE}，它们通过引入额外的编码器（如自动编码器 Autoencoder）来将视觉或其他模态的信息融入实体的向量表示中。最新的研究（如 \textbf{MMKG-T5}）则利用强大的预训练Transformer模型，实现了对结构、文本、视觉等多模态信息更深层次的融合。
\end{itemize}
这些向量化后的知识，可以直接服务于认知校准框架中的“音系特征编码”环节，以及自适应学习算法中的内容推荐模块，真正实现了从结构化知识到AI模型可利用特征的转化，其概览如图 \ref{fig:mmkg_pipeline} 所示

\begin{figure}[H]
    \centering
    \includegraphics[width=0.9\textwidth]{多模态知识图谱概览.pdf}
    \caption{多模态知识图谱概览}
    \label{fig:mmkg_pipeline}
\end{figure}

\section{知识图谱的评估与应用}
构建完成的“多模态方言声学知识图谱”不仅需要在理论设计上科学合理，更需在实践中证明其自身的结构质量与数据准确性。同时，它作为本项目所有上层算法的基石，其核心价值最终体现在对这些算法的赋能上。因此，本节将首先对其进行严格的内在质量评估，然后阐述其在后续核心算法中的具体应用方式。

\subsection{内在质量评估}
为保证知识图谱作为一个高质量静态数据制品的可靠性，本项目仅采用内在评估（Intrinsic Evaluation）的方式，从结构和知识表示两个层面，对图谱自身的质量进行检验。

\paragraph{结构完整性与一致性评估}
在将数据正式载入图数据库之前，本项目首先对图谱的结构和模式（Schema）进行了严格的验证。第一，本项目采用标准的本体论推理机（Ontology Reasoner），对本项目扩展后的OntoLex-Lemon模型进行了逻辑一致性检查，确保其中不存在任何矛盾的公理或无法被满足的类定义。第二，本项目对所有生成的三元组进行了类型约束验证，即检查每一个关系（Property）的主语（Domain）和宾语（Range）是否符合本项目在Schema中定义的实体类型。例如，确保`hasAudioExample`这一关系的宾语必须指向一个`Multimodal Asset`类型的实体。通过这些验证，本项目从根本上保证了知识图谱的结构严谨性与数据一致性。

\paragraph{知识嵌入质量评估：链接预测}
为评估知识图谱中存储的关系是否具有规律性、是否能被机器学习模型有效理解和表示，本项目采用了标准的\textbf{链接预测（Link Prediction）}任务。此任务旨在检验知识图谱嵌入模型的表示能力。本项目评估并比较了多种嵌入模型，从传统的TransE、RotatE到为本项目设计的多模态嵌入模型。具体而言，本项目随机地从图谱中移除一部分三元组的头实体 `(?, r, t)` 或尾实体 `(h, r, ?)`，让模型根据学习到的实体和关系向量来预测缺失的部分。

评估指标采用 `Hits@10`（预测结果排名前10的命中率）和 `MRR`（平均倒数排名）。实验结果呈现出两个关键结论：首先，所有模型均取得了较高的分数，这证明了本知识图谱的整体结构是清晰、规整且包含丰富逻辑关联的，能够被嵌入模型有效学习。其次，多模态嵌入模型在预测涉及多媒体资产的链接（例如，`(音频片段A, hasSpectrogram, ?)`）时，其性能显著优于传统模型。这一结果不仅验证了本项目嵌入算法的有效性，更从数据层面雄辩地证明了，将多模态特征直接融入知识表示，对于构建一个完备的声学知识图谱是至关重要的。

\subsection{在核心算法中的应用}
内在评估证实了本知识图谱作为一个静态数据制品的高质量。然而，其最终的研究价值，体现在它如何作为核心知识源，去驱动和赋能本项目后续提出的高级AI模型。本节旨在阐述其具体的应用方式，而关于这些应用所带来的实际性能提升的外在评估，本项目将在后续的实验章节中进行集中和详细的论证。

\paragraph{作为“文脉智能体”的结构化长时记忆}
“文脉智能体”框架虽然以大语言模型为核心，但LLM本身存在“幻觉”和知识更新不及时的问题。本知识图谱在此扮演了智能体系统的、可信赖的“\textbf{结构化长时记忆}”。当处理专业问题时，声韵智能体能够对知识图谱进行精准的图查询，获取可验证的、精确的语言学事实。这种基于结构化知识的检索，相比于LLM自身模糊的、概率性的参数化记忆，在准确性和可解释性上具有压倒性优势，是确保“文脉智能体”专业性的根本保障。

\paragraph{作为“自适应教学模型”的内容与策略基座}
知识图谱的精细化设计，是整个自适应教学模型得以运转的“内容与策略基座”。首先，图谱中细粒度的知识节点使得“\textbf{学习者状态向量}”$S_u$的构建能够精准到具体的知识点上。其次，图谱中丰富的\textbf{多模态资产}，是生成“个性化教学包”的“弹药库”，是实现本项目设计的、数据驱动的、声学导向教学范式的物质基础。

\paragraph{作为“自进化模型”的真值参照与质量标尺}
在“音系自进化生成与校准模型”中，知识图谱扮演了“\textbf{真值参照}”与“\textbf{质量标尺”}的双重角色。模型进行发音诊断时所检索到的标准示范$z^*$，正是源自知识图谱中权威的、高质量的资产。同时，它也为“教师模型”过滤劣质生成数据提供了黄金标准，确保了模型在持续的迭代中，始终朝着更专业、更准确的方向演进。
